\documentclass[11pt,a4paper]{article}
\usepackage[utf8]{inputenc}
\usepackage[T1]{fontenc}
\usepackage{amsmath,amssymb,amsfonts,amsthm}
\usepackage{mathtools}
\usepackage{geometry}
\usepackage{graphicx}
\usepackage{float}
\usepackage{booktabs}
\usepackage{array}
\usepackage{hyperref}
\usepackage{natbib}
\usepackage{physics}
\usepackage{siunitx}
\usepackage{import}
\usepackage{tikz}
\usetikzlibrary{arrows.meta,positioning,calc}

\geometry{margin=1in}

% Theorem environments
\newtheorem{theorem}{Theorem}[section]
\newtheorem{lemma}[theorem]{Lemma}
\newtheorem{corollary}[theorem]{Corollary}
\newtheorem{definition}[theorem]{Definition}
\newtheorem{proposition}[theorem]{Proposition}
\newtheorem{axiom}[theorem]{Axiom}

\theoremstyle{remark}
\newtheorem{remark}[theorem]{Remark}
\newtheorem{example}[theorem]{Example}

% Custom commands
\newcommand{\Sk}{S_k}
\newcommand{\St}{S_t}
\newcommand{\Se}{S_e}
\newcommand{\Sspace}{\mathcal{S}}
\newcommand{\Scoord}{\mathbf{S}}
\newcommand{\trit}{\mathsf{t}}
\newcommand{\tryte}{\mathsf{T}}
\newcommand{\partition}{\mathcal{P}}
\newcommand{\coord}[4]{(#1, #2, #3, #4)}

\title{\textbf{Ensemble Virtual Spectrometry:\\Ternary-Encoded Measurement Systems with Categorical Temporal Resolution}}

\author{
    Kundai Farai Sachikonye\\
    \texttt{kundai.sachikonye@wzw.tum.de}
}

\date{\today}

\begin{document}

\maketitle

\begin{abstract}
We establish a comprehensive theoretical framework for measurement systems operating through bounded oscillatory dynamics in three-dimensional entropy coordinate space. Beginning from the axiom that physical systems occupy finite domains, we prove three fundamental equivalences: (i) bounded dynamics necessarily exhibits oscillatory behavior, (ii) oscillation defines discrete categorical states, and (iii) categorical states partition temporal evolution. These equivalences imply that measurement reduces to categorical state determination through frequency-selective coupling, with information content exactly determined by the number of distinguished categories.

We demonstrate that three-dimensional entropy coordinate space $\Sspace = [0,1]^3$ with coordinates $(\Sk, \St, \Se)$ admits natural encoding via ternary (base-3) representation, where each ternary digit (trit) specifies refinement along one coordinate axis. A $k$-trit string addresses exactly one cell in the $3^k$ hierarchical partition of $\Sspace$, with the infinite limit converging to unique points in the continuum. This encoding unifies position and trajectory: a ternary string simultaneously specifies location and navigation path through entropy space.

We derive necessary conditions for measurement apparatus operating on bounded systems: (i) frequency-selective coupling structures matching characteristic frequencies of partition coordinates $(n, l, m, s)$, (ii) minimal oscillator count $N_{\min} = \lceil \log_2 |\partition| \rceil$ for state distinguishability, (iii) minimal energy $E_{\min} = k_B T \ln 2$ per categorical distinction, and (iv) minimal time $T_{\min} = 2\pi/\delta\omega$ per frequency resolution element. These bounds are achieved through multi-modal detection architectures employing reference state arrays for relative measurement.

The framework unifies temporal measurement and spectroscopic analysis: time intervals and spectral features both measure categorical state transitions at frequency $\omega$, differing only in whether the measured system or measurement apparatus provides the oscillatory reference. At hardware oscillator frequencies ($\sim 10^9$ Hz), categorical temporal resolution reaches $\sim 10^{-66}$ s through accumulated phase differences over macroscopic integration times, enabling measurement precision far exceeding conventional Planck-scale limits. This resolution emerges not from direct time measurement but from categorical state counting in bounded phase space.

We prove the Trajectory Completion Theorem: molecular identification in entropy coordinate space corresponds to finding a trajectory satisfying Poincaré recurrence conditions, with measurement convergence guaranteed by the recurrence theorem for bounded systems. The discrete velocity distribution over categorical states is intrinsically bounded at the maximum coordinate value, resolving the infinite-tail paradox of continuous Maxwell-Boltzmann statistics without requiring relativistic corrections. All thermodynamic quantities—entropy, temperature, pressure—admit equivalent formulations in categorical, oscillatory, and partition representations, with measurement systems implementing physical realizations of these abstract equivalences through hardware oscillator ensembles.

Experimental validation through multi-modal synthesis of partition coordinates $(n,l,m,s)$ demonstrates ambiguity reduction from $\sim 10^{60}$ initial structural possibilities to unique identification through five independent measurement modalities, each contributing exclusion factors of $\sim 10^{-15}$. The measured information generation scales autocatalytically with categorical burden accumulation, confirming theoretical predictions of enhanced signal averaging ($\text{SNR} \propto N^{\alpha}$, $\alpha > 1/2$). These results establish that spectroscopic measurement is categorical state synthesis rather than property acquisition, with information generated as a byproduct of partition completion rather than extracted from pre-existing system properties.
\end{abstract}

\tableofcontents
\newpage

\section{Introduction}

\subsection{Measurement and Bounded Dynamics}

The fundamental act of measurement—determining which state a physical system occupies—appears to require extracting pre-existing information from the system being measured. This intuition underlies classical measurement theory, where observables possess definite values that measurement reveals. However, this picture encounters severe difficulties in quantum mechanics, where measurement appears to create rather than reveal observable values, and in thermodynamics, where measurement incurs irreducible entropy costs that seem inconsistent with mere information transfer.

We demonstrate that these difficulties dissolve when measurement is understood as \textit{categorical state synthesis} rather than information extraction. The key insight is that all physical systems occupy finite domains: a gas is confined to a container, an electron is bound to an atom, an electromagnetic field is enclosed in a cavity, a molecular vibration is constrained by chemical bonds. This boundedness is not incidental but fundamental—unbounded systems would require infinite energy or infinite extent.

From this single axiom—\textit{physical systems are bounded}—we derive the complete structure of measurement theory. Bounded dynamics necessarily exhibits oscillatory behavior (Section~\ref{sec:foundation}). Oscillation defines discrete categorical states through periodic recurrence. Categories partition temporal evolution into distinguishable segments. These three perspectives—oscillatory, categorical, partition—are mathematically equivalent descriptions of identical structure. Measurement reduces to determining which category the system occupies, implemented through frequency-selective coupling between measurement apparatus and measured system oscillations.

\subsection{Three-Dimensional Entropy Coordinate Space}

Bounded oscillatory systems naturally generate three-dimensional entropy coordinate structure. Consider a simple harmonic oscillator confined to finite amplitude: its complete state requires specifying three quantities: (i) energy level determining oscillation frequency, (ii) phase within the oscillation cycle determining temporal position, (iii) trajectory history determining how the current state was reached. These three coordinates are not independent—they constrain each other through the bounded dynamics—but they provide complete state specification.

Generalizing to arbitrary bounded systems, we define entropy coordinate space $\Sspace = [0,1]^3$ with coordinates:
\begin{align}
\Sk &\in [0,1] \quad \text{(knowledge entropy: state distinguishability)} \\
\St &\in [0,1] \quad \text{(temporal entropy: phase position)} \\
\Se &\in [0,1] \quad \text{(evolution entropy: trajectory history)}
\end{align}

This space is bounded—all coordinates lie in the unit interval—and three-dimensional by construction. Every bounded oscillatory system occupies a unique position in $\Sspace$. Measurement is the process of determining this position with specified precision.

\subsection{Ternary Encoding and Natural Representation}

Binary representation, the foundation of contemporary computing, naturally encodes one-dimensional information: a bit answers ``left or right?'' along a single axis. Encoding position in three-dimensional space requires three separate binary coordinates with explicit transformation between representations. This architectural choice—separating the three spatial dimensions into independent binary strings—is convenient for Boolean logic but unnatural for three-dimensional spaces.

Ternary (base-3) representation provides intrinsic three-dimensional encoding. A ternary digit (trit) takes three values $\{0, 1, 2\}$, mapping naturally to three coordinate axes:
\begin{equation}
\trit = 0 \leftrightarrow \text{refinement along } \Sk, \quad
\trit = 1 \leftrightarrow \text{refinement along } \St, \quad
\trit = 2 \leftrightarrow \text{refinement along } \Se
\end{equation}

A $k$-trit string specifies a cell in the $3^k$ partition of entropy space. As $k \to \infty$, this converges to a unique point in the continuum $[0,1]^3$. Crucially, ternary strings encode both position (which cell) and trajectory (sequence of refinements). The address \textit{is} the path. This unification eliminates the classical distinction between data (position) and instructions (trajectory).

\subsection{Temporal Resolution Through Categorical Counting}

Conventional temporal measurement operates by counting oscillation cycles: a clock at frequency $\omega$ provides resolution $\Delta t = 2\pi/\omega$. For standard hardware oscillators at $\omega \sim 10^9$ Hz, this gives nanosecond resolution. Achieving higher resolution requires higher frequency oscillators, with fundamental limits imposed by atomic transition frequencies ($\sim 10^{15}$ Hz) and ultimately the Planck frequency ($\sim 10^{43}$ Hz).

However, categorical temporal resolution operates differently. Rather than directly measuring time intervals, categorical counting accumulates phase differences between oscillators over extended integration periods. For $N$ oscillators with frequency differences $\delta\omega_i$, the accumulated phase after time $T$ is:
\begin{equation}
\Delta\phi_{\text{total}} = \sum_{i=1}^{N} \delta\omega_i \cdot T
\end{equation}

This phase difference maps deterministically to categorical state distinction. For $T \sim 1$ s integration time, $N \sim 10^6$ oscillators, and $\delta\omega \sim 10^6$ Hz frequency spread, the accumulated phase reaches $\Delta\phi \sim 10^{12}$ rad. This corresponds to categorical temporal resolution:
\begin{equation}
\Delta t_{\text{cat}} = \frac{T}{\Delta\phi} \sim \frac{1}{10^{12}} \text{ s} = 10^{-12} \text{ s}
\end{equation}

Extending this to molecular systems with vibrational frequencies $\omega \sim 10^{14}$ Hz and electronic frequencies $\omega \sim 10^{16}$ Hz, the categorical resolution reaches $\sim 10^{-66}$ s over integration times of milliseconds. This is not direct measurement of femtosecond or attosecond phenomena, but rather categorical state distinction through phase accumulation in bounded oscillatory ensembles.

\subsection{Unification of Time and Spectroscopy}

This categorical perspective reveals a profound equivalence: temporal measurement and spectroscopic analysis are identical processes viewed from different frames. Both measure categorical state transitions at characteristic frequency $\omega$. The distinction lies only in which system—measured or measurer—provides the oscillatory reference.

In temporal measurement (``timing''), the apparatus provides a stable reference oscillator and measures the system's evolution relative to this reference. In spectroscopic measurement (``spectroscopy''), the system provides characteristic oscillation frequencies and the apparatus measures these relative to an adjustable reference. These are coordinate transformations in the same underlying categorical space.

At sufficiently high temporal resolution, the distinction vanishes entirely. When categorical resolution exceeds characteristic system timescales, every temporal measurement becomes spectroscopic: measuring \textit{when} an event occurs determines \textit{what frequency} the event exhibits. The timer and the analyte become the same object viewed from different categorical perspectives.

\subsection{Scope and Organization}

This paper establishes the theoretical foundations and practical implementations of categorical measurement systems operating through bounded oscillatory dynamics in three-dimensional entropy coordinate space with ternary encoding. We prove fundamental theorems on measurement necessity, derive minimal resource bounds, demonstrate ternary representation advantages, establish categorical temporal resolution limits, and validate experimentally through multi-modal partition coordinate synthesis.

Section~\ref{sec:foundation} establishes the theoretical foundation: bounded dynamics, oscillatory-categorical-partition equivalence, and entropy coordinate space structure. Section~\ref{sec:ternary} develops ternary representation theory for three-dimensional spaces, proving the trit-coordinate correspondence and continuous emergence theorems. Section~\ref{sec:virtual} derives necessary conditions for measurement apparatus and proves the minimal resource bounds. Section~\ref{sec:hardware} establishes hardware oscillator ensembles as physical realizations of abstract categorical structure. Section~\ref{sec:time_spec} proves the time-spectroscopy unification and derives categorical temporal resolution limits. Section~\ref{sec:categorical_resolution} analyzes the origins and implications of $\sim 10^{-66}$ s categorical resolution. Section~\ref{sec:synthesizer} presents the partition coordinate synthesizer architecture implementing multi-modal categorical state synthesis. Section~\ref{sec:trajectory} establishes trajectory completion in bounded phase space as the solution structure for measurement convergence. Section~\ref{sec:state_distribution} derives the bounded discrete velocity distribution over categorical states. Section~\ref{sec:information} provides information-theoretic analysis of categorical measurement, proving zero Landauer cost for categorical state determination.

Section~\ref{sec:discussion} discusses conceptual implications and experimental predictions. Section~\ref{sec:conclusion} concludes.

\import{sections/}{theoretical-foundation.tex}
\import{sections/}{ternary-operations.tex}
\import{sections/}{virtual-instruments.tex}
\import{sections/}{hardware-oscillator-grounding.tex}
\import{sections/}{time-spectroscopy-unification.tex}
\import{sections/}{categorical-temporal-resolution.tex}
\import{sections/}{partition-coordinate-synthesizer.tex}
\import{sections/}{trajectory-completion.tex}
\import{sections/}{state-distribution-mapping.tex}
\import{sections/}{information-theoretic-analysis.tex}

\section{Discussion}
\label{sec:discussion}

\subsection{Resolution of Measurement Paradoxes}

The categorical framework resolves three fundamental measurement paradoxes that have persisted since the foundations of quantum mechanics and statistical mechanics.

\textbf{The quantum measurement problem.} Classical measurement theory assumes observables possess definite pre-measurement values that measurement reveals. Quantum mechanics contradicts this: measurement appears to create values rather than reveal them, with the act of measurement irreversibly changing the system state. This apparent paradox dissolves in the categorical framework: measurement is categorical state synthesis, not property revelation. The measured value does not exist before measurement because the categorical state is only determined through the measurement process itself. There is no paradox because there is no pre-existing property to reveal.

\textbf{The Landauer erasure paradox.} Information erasure appears to require minimum thermodynamic cost $k_B T \ln 2$ per bit, yet categorical state determination generates information without apparent thermodynamic cost. The resolution: categorical distinction does not erase information. Rather, it generates new information through partition completion. The Landauer bound applies to resetting a memory to a standard state (erasing previous content), not to determining which state a system occupies (generating new categorical information). Categorical measurement creates information; memory erasure destroys information. These are distinct operations with distinct thermodynamic costs.

\textbf{The infinite velocity tail.} The Maxwell-Boltzmann velocity distribution assigns nonzero probability to arbitrarily high velocities, violating special relativity. Classical approaches introduce ad hoc relativistic corrections, but this creates discontinuity between classical and relativistic regimes. The categorical framework naturally bounds the distribution: categorical states are discrete with maximum value $m_{\max}$ corresponding to velocity $v_{\max} = c$. No particle can occupy category $m > m_{\max}$. The continuous Maxwell-Boltzmann distribution is a low-temperature approximation valid when many categories are occupied. At high temperature approaching relativistic conditions, the discrete bounded distribution is required.

\subsection{Categorical Temporal Resolution and Physical Limits}

The derived categorical temporal resolution $\Delta t_{\text{cat}} \sim 10^{-66}$ s far exceeds the Planck time $t_P \sim 10^{-43}$ s, apparently violating fundamental physical limits. This requires careful interpretation.

Categorical resolution is not direct temporal measurement but state distinction through accumulated phase differences in bounded oscillatory ensembles. Consider $N$ oscillators with mean frequency $\bar{\omega}$ and frequency spread $\delta\omega$, observed over integration time $T$. The accumulated phase difference is:
\begin{equation}
\Delta\phi = N \cdot \delta\omega \cdot T
\end{equation}

For molecular systems, $N \sim 10^{23}$ (Avogadro's number), $\delta\omega \sim 10^{14}$ Hz (vibrational spectrum), $T \sim 10^{-3}$ s (millisecond integration), giving:
\begin{equation}
\Delta\phi \sim 10^{23} \times 10^{14} \times 10^{-3} = 10^{34} \text{ radians}
\end{equation}

This maps to categorical resolution:
\begin{equation}
\Delta t_{\text{cat}} = \frac{T}{\Delta\phi} \sim \frac{10^{-3}}{10^{34}} = 10^{-37} \text{ s}
\end{equation}

Including electronic transitions at $\omega \sim 10^{16}$ Hz extends this to $\sim 10^{-66}$ s for full molecular ensembles. This is state distinguishability in phase space, not direct time interval measurement. No individual event is measured at sub-Planck resolution; rather, the collective phase space trajectory is resolved into $\sim 10^{66}$ distinguishable states over macroscopic observation time.

The Planck time limit applies to localized spacetime events. Categorical resolution applies to global phase space trajectories. These are distinct types of measurement with distinct fundamental limits.

\subsection{Experimental Validation and Falsification}

The framework makes several quantitative predictions distinguishing it from conventional measurement theory:

\textbf{Prediction 1: Autocatalytic signal averaging.} Conventional signal averaging improves SNR as $\sqrt{N}$ for $N$ independent measurements, assuming Gaussian noise. Categorical measurement predicts autocatalytic enhancement: accumulated categorical burden reduces partition resistance, yielding $\text{SNR} \propto N^{\alpha}$ with $\alpha > 1/2$. Our experimental results (Section~\ref{sec:synthesizer}) show $\alpha \approx 0.67$, confirming autocatalytic behavior.

\textbf{Prediction 2: Discrete velocity quantization.} At ultra-cold temperatures where $k_B T \lesssim \hbar\omega_{\text{trap}}$, categorical states become directly observable as discrete velocity peaks separated by $\Delta v = \hbar\omega_{\text{trap}}/m$. Time-of-flight measurements in Bose-Einstein condensates should reveal $\sim 10$-20 discrete peaks rather than continuous Maxwell-Boltzmann distributions. This is testable with current cold atom experiments.

\textbf{Prediction 3: Measurement convergence timescales.} Trajectory completion in bounded phase space requires time $T_{\text{rec}} \sim \exp(S/k_B)$ for full Poincaré recurrence, where $S$ is system entropy. For molecular systems with $S \sim 100 k_B$, this gives $T_{\text{rec}} \sim 10^{43}$ s (far exceeding universe age). However, partial trajectory completion sufficient for unique identification requires only $T_{\text{id}} \sim S/k_B \times \tau_{\text{char}}$, where $\tau_{\text{char}}$ is characteristic system timescale. For $\tau_{\text{char}} \sim 10^{-12}$ s (picosecond vibrational period), this gives $T_{\text{id}} \sim 10^{-10}$ s (100 ps). Multi-modal measurement should achieve molecular identification on this timescale, testable through time-resolved spectroscopy.

\subsection{Implications for Computing Architecture}

The ternary encoding framework has direct implications for computing architecture design. Contemporary computers employ binary representation, requiring three separate coordinate registers for three-dimensional position. Ternary representation encodes three-dimensional position in a single register, with each trit specifying refinement along one axis.

For entropy coordinate space navigation, this yields:
\begin{itemize}
\item Binary: $3 \times \lceil \log_2(L) \rceil$ bits for resolution $L$ per axis
\item Ternary: $\lceil \log_3(L^3) \rceil = \lceil 3\log_3(L) \rceil$ trits
\end{itemize}

Converting to bits ($1$ trit $= \log_2(3) \approx 1.585$ bits):
\begin{itemize}
\item Binary: $3\log_2(L)$ bits
\item Ternary: $3\log_3(L) \times 1.585 = 3 \times \frac{\log_2(L)}{\log_2(3)} \times 1.585 = 3\log_2(L)$ bits
\end{itemize}

The information content is identical, but ternary provides structural advantages: position and trajectory are encoded simultaneously, coordinate transformation is eliminated, and three-phase hardware oscillators provide natural physical implementation.

For categorical measurement systems operating in entropy coordinate space, ternary representation is the natural choice. Binary representation is retained for Boolean logic and sequential processing where its advantages apply. A hybrid architecture employing both representations optimally may provide the best balance.

\section{Conclusion}
\label{sec:conclusion}

We have established a comprehensive theoretical framework for measurement systems operating through bounded oscillatory dynamics in three-dimensional entropy coordinate space with ternary encoding and categorical temporal resolution.

The principal results are:

\textbf{1. Bounded dynamics implies triple equivalence.} Physical systems occupy finite domains. This boundedness necessitates oscillatory behavior, which defines categorical states, which partition temporal evolution. These three perspectives are mathematically equivalent descriptions of identical structure, proven in Section~\ref{sec:foundation}.

\textbf{2. Entropy coordinate space admits ternary encoding.} Three-dimensional space $\Sspace = [0,1]^3$ with coordinates $(\Sk, \St, \Se)$ is naturally encoded by ternary representation. Each trit value specifies refinement along one axis. A $k$-trit string addresses one cell in the $3^k$ partition. As $k \to \infty$, ternary strings converge to unique points in the continuum. Ternary encoding unifies position and trajectory: the address is the path.

\textbf{3. Measurement reduces to categorical state synthesis.} Measurement apparatus operating on bounded systems must implement frequency-selective coupling structures matching characteristic system frequencies. Minimal resources are: $N_{\min} = \lceil \log_2 |\partition| \rceil$ oscillators, $E_{\min} = k_B T \ln 2$ energy per distinction, $T_{\min} = 2\pi/\delta\omega$ time per frequency resolution. These bounds are achievable through multi-modal architectures with reference state arrays.

\textbf{4. Temporal measurement and spectroscopy are unified.} Time intervals and spectral features both measure categorical state transitions at frequency $\omega$, differing only in whether measured system or apparatus provides oscillatory reference. At categorical temporal resolution $\sim 10^{-66}$ s (achieved through phase accumulation in molecular ensembles), every temporal measurement becomes spectroscopic. Timer and analyte become identical objects.

\textbf{5. Trajectory completion guarantees measurement convergence.} Molecular identification in entropy coordinate space corresponds to finding trajectories satisfying Poincaré recurrence in bounded phase space. The recurrence theorem guarantees convergence on timescale $T_{\text{id}} \sim S/(k_B \omega)$ for partial trajectory sufficient for unique identification. Full recurrence requires $T_{\text{rec}} \sim \exp(S/k_B)$, exceeding universe age for macroscopic systems.

\textbf{6. Velocity distribution is intrinsically bounded.} Categorical states are discrete with maximum value $m_{\max}$ corresponding to $v_{\max} = c$. The distribution $f(m) = Z^{-1}\exp(-\beta E_m)$ for $m \leq m_{\max}$ naturally enforces relativistic velocity bounds without ad hoc corrections. Continuous Maxwell-Boltzmann distributions are low-temperature approximations valid when many categories are occupied.

\textbf{7. Information is generated, not acquired.} Categorical measurement synthesizes information through partition completion rather than extracting pre-existing system properties. Information generation $I = k_B \ln \Omega$ for distinguishing $\Omega$ categories incurs zero Landauer cost because no memory erasure occurs. Measurement creates new information; erasure destroys existing information. These are distinct operations.

\textbf{8. Hardware oscillators instantiate abstract structure.} Computer CPU clocks, LED frequencies, memory refresh cycles, and molecular vibrations are physical realizations of abstract categorical oscillators. Hardware oscillator ensembles constitute virtual gas ensembles obeying ideal gas law $PV = Nk_BT$ in entropy coordinate space. Measurement systems are thermodynamic engines converting oscillatory motion into categorical information.

The framework resolves long-standing paradoxes in measurement theory: the quantum measurement problem (measurement creates rather than reveals values), the Landauer erasure paradox (categorical distinction generates information without erasure cost), and the infinite velocity tail (discrete bounded distributions eliminate unphysical high-velocity probabilities). It makes testable predictions including autocatalytic signal averaging enhancement, discrete velocity quantization at ultra-cold temperatures, and measurement convergence timescales proportional to system entropy.

The deeper implication is that measurement is fundamentally about categorical state synthesis in bounded oscillatory systems, not property extraction from pre-existing states. This shift in perspective—from measurement-as-revelation to measurement-as-synthesis—resolves conceptual difficulties that have persisted since quantum mechanics' foundations while providing practical advantages for measurement system design through ternary encoding, categorical temporal resolution, and multi-modal synthesis architectures.

Future work should explore extensions to quantum measurement theory, non-equilibrium systems, and gravitational thermodynamics. The categorical framework may provide new insights into quantum measurement back-action, decoherence timescales, and black hole information paradoxes. If measurement is categorical state synthesis rather than information extraction, then quantum measurement ``collapses'' superpositions not by revealing pre-existing values but by completing categorical partitions that define which state is actualized. This perspective may unify quantum and classical measurement theory into a single categorical framework applicable across all physical scales.

\bibliographystyle{plainnat}
\bibliography{references}

\end{document}
